\addcontentsline{toc}{chapter}{Abbreviations and Acronyms}
\chapter*{Abbreviations and Acronyms}

% The longtable environment should break the table properly to multiple pages, 
% if needed

\noindent
%\begin{longtable}{@{}p{0.25\textwidth}p{0.7\textwidth}@{}}
\begin{longtable}{@{}p{0.19\textwidth}p{0.8\textwidth}@{}}
IoT & Internet of Things \\
API  & Application Programming Interface \\
SO   & Smart Object \\
HTML & HyperText Markup Language; The language of website content. \\
JS   & JavaScript; A programming language. Most significantly, browsers only allow JS for website client-side functionality. \\
CSS  & Cascading Style Sheets; The language of website styling. \\
HTTP & Hypertext Transfer Protocol; The core protocol of the Web. Used to transfer HTML and other files. \\
O-MI & Open Messaging Interface; An IoT protocol standard \\
O-DF & Open Data Format; An IoT data standard \\
P2P & Peer to Peer; A term for communications directly from peer to another peer, instead of via a dedicated server. \\
RSOS & Reconfigurable Smart Object System; The system architecture proposed in this study. \\
PnP & Plug and Play; A term used for easy-to-install peripherals, without user intervention. \\
%AWS & Amazon Web Services; The most well known cloud platform. \\
WLAN & Wireless Local Area Network; See LAN. \\
LAN & Local Area Network; A computer network that spans a limited physical area, such as one apartment or building. \\
Wi-Fi & A family of wireless network protocols, often used for WLAN. \\
pub/sub & Publish-subscribe pattern, where participants can either publish new information or subscribe to consume it. \\
MQTT & Message Queuing Telemetry Transport; A lightweight publish-subscribe protocol. \\
CoAP & Constrained Application Protocol; An application data transfer protocol intended for constrained devices.\\
REST & Representational State Transfer; A software architecture style, simplifying HTTP services. \\
%, which simplifies web application development and improves interoperability between systems. \\
TCP & Transmission Control Protocol; A fundamental Internet protocol to provide reliable and ordered stream of data. \\
UDP & User Datagram Protocol; An Internet protocol to send messages without ordering or delivery guarantees. \\
UPnP & Universal Plug and Play; A group of protocols that allow networked devices to find and establish connections to each other automatically in a LAN. \\
DNS & Domain Name System; A system and protocol that is mainly used to retrieve IP addresses for hostnames. \\
mDNS & Multicast DNS; A modified DNS protocol; Allows DNS queries in a LAN without a name server with multicast UDP. \\
DNS-SD & DNS Service Discovery; mDNS companion standard to allow multiple results to find all hosts providing a specified service. \\
LLMNR & Link-Local Multicast Name Resolution; Similar to mDNS. \\
CRUD & Create, Read, Update and Delete; A term for the required basic operations to create a persistent data storage. \\
SOAP & Simple Object Access Protocol; Allows transfer of structured information and invoking processes on a server. \\
XML & Extensible Markup Language; A document markup format that is both human-readable and machine-readable. \\
TTL & Time-To-Live; A parameter in O-MI that defines the maximum amount of time the request should be answered in. \\
%USB & Universal Serial Bus; An industry standard for computer peripheral device communication. \\
URL & Uniform Resource Locator; Also known as a web address; A format to specify a location on a computer network. \\
WADL & Web Application Description Language; An XML format to describe HTTP-based web services. \\
IP & Used to refer either the Internet Protocol or a numerical address used in it. \\
DPWS & Devices Profile for Web Services; Similar standard to UPnP, but aimed more towards web service technologies. \\
WS-Discovery & Web Services Dynamic Discovery; A multicast service discovery protocol for local networks. \\
VLAN & Virtual Local Area Network; A standard that allows creation of many virtual LANs in one physical LAN. \\
SoC & System on Chip; An integrated circuit that contains most components of a computer system, instead of separate parts. \\
ESP32-S2 & The commercial name of the SoC used in this study. \\
RAM & Random Access Memory; A memory type, which allows access to any of its memory locations in roughly the same time. \\
SRAM & Static Random Access Memory; A type of semiconductor RAM that requires no refreshing circuitry to retain the data. \\
PSRAM & PseudoStatic Random Access Memory; Similar to static RAM in interface but handles memory refreshing internally. \\
UART & Universal Asynchronous Receiver-Transmitter; A device for configurable serial communication between devices. \\
RGB LED & Red Green Blue Light Emitting Diode; A light that has configurable color channels to achieve any perceived light color. \\
GPIO & General-Purpose Input/Output; A hardware digital signal pin on an integrated circuit, that can be configured to various input or output operations by the firmware. \\
CPU & Central Processing Unit; An integrated circuit that performs operations, as specified by instructions of a program. \\
ROM & Read-Only Memory; A type of memory that cannot be written. Provides commonly used program code or data in an integrated circuit to save physical space or faster speed. \\
RTC & Real-Time Clock; An electronic device that retains and measures passage of clock time and date. \\
IDE & Integrated Development Environment; A software that combines software development tools to one integrated package. \\
OTA & Over-The-Air; A term for installing firmware or software updates remotely, usually over a wireless connection. \\
%UI & User Interface; Used to operate a software or device. \\
RDF & Resource Description Framework; A family of specifications used to describe resources in machine-readable way. \\
app & Short for application; Typically refers to a smartphone app. \\
OS & Operating System; A system software that manages hardware resources and running of programs. \\
NAT & Network Address Translation; A method of changing an IP address in network packets when routing between networks. \\% Common in routers between private networks and Internet. \\
WPS & Wi-Fi Protected Setup; A standard that attempts to automate wireless network connection setup securely. \\
NFC & Near-Field Communication; A group of wireless communication protocols, intended to restrict range to a few cm. \\
eOMI & Embedded O-MI; The embedded system implementation of O-DF/O-MI in this study. \\
UTF & Unicode Transformation Format; A Format to encode text characters of most of the world's writing systems. \\
DTD & Document Type Definition; Declaration of a document type used in languages such as XML. \\
MCU & MicroController Unit; A small computer in an integrated circuit. \\
NTP & Network Time Protocol; A network protocol for clock synchronization between computer systems. \\


\end{longtable}
% EXAMPLES
%2k/4k/8k mode & COFDM operation modes \\
%3GPP & 3rd Generation Partnership Project \\ 
%ESP & Encapsulating Security Payload; An IPsec security protocol \\ 
%FLUTE  & The File Delivery over Unidirectional Transport protocol \\ 
%e.g.& for example (do not list here this kind of common acronymbs or abbreviations, but only those that are essential for understanding the content of your thesis. \\ 
%note & Note also, that this list is not compulsory, and should be omitted if you have only few abbreviations
